% Transcription — fonds Alexandre Grothendieck, shelfmark 137, batch 5 (pages 81-85).
% Established from the facsimile archives/batches/137/batch-05.pdf (a copy of
% 137.pdf fetched through the project's relay,
% https://grothendieck-relay.onrender.com/source/137.pdf; PDF page k is
% archivists' page 80+k, no cover sheet), per the skill transcribe-grothendieck.
% This is the last of the folder's five batches (85 pages: 1-20, 21-40, 41-60,
% 61-80, 81-85), transcribed in parallel; none of batches 1-4 existed when this
% one was written, so the notation below was fixed from pages 79-80 of the
% facsimile (batch 4), which these pages continue.
%
% Pass: Opus 5.5 (claude-opus-5-5), 2026-09-26 — first pass, unchecked against the pages by a human.
%
% Pages skipped: none. Pages summarised: none. All five pages are his
% mathematics and are transcribed.
%
% His pagination: none visible on these five leaves; no numbered run begins
% or ends here. His labels carried over from page 80 are the conditions (*)
% (associativity of the alpha_{p,q}), (**) and (***) (unit conditions, stated
% on p. 81); the axioms « a) à f) » cited on p. 83 are not on these leaves.
% Pages 81-82 are in blue ballpoint and pencil (blue retracings over pencil),
% pages 82 and 85 on music-staff paper; pages 83-85 in black ink, written
% faster and more legibly. The last line of p. 82 is in black ink, in another
% stroke; it is given as a \marginal with its words mostly \ill{}.
%
% What the batch is about. The end of the construction of « catégories
% polynomiales » from a ⊗-additive category P_1 with divided-power-like
% functors. P. 81: conversely, given P_1 ⊗-additive (ACU), with a ⊗-functor
% Γ^* : (P_1, ⊕) -> (P_1^grad, ⊗), Γ^*(0) = 1, Γ^*(X ⊕ Y) ≃ Γ^*(X) ⊗ Γ^*(Y),
% Γ^0 = 1, Γ^1 = id, and maps α_{p,q} : Γ^{pq} -> Γ^p ∘ Γ^q satisfying the
% associativity condition of p. 80 and the unit conditions (**), (***). P. 82:
% one defines P_p(X,Y) = Hom(Γ^p X, Y) and the composition
% P_p(X,Y) × P_q(Y,Z) -> P_{pq}(X,Z), g∘f = g Γ^q(f) α^X_{q,p}; (*) gives
% associativity, (**) the compatibility with bifunctoriality, (***) that
% g∘f = g for g of degree 0 — plus the axiom 1 ≃ Γ^q(1), which he says he had
% forgotten among the preliminary axioms. P. 83: for X = ⊕ X_i, the
% decomposition of Γ^n(X) gives that of P_n(∏X_i, Y) into the
% P_{ν_1...ν_r}; « on vérifie alors les axiomes a) à f) »; a polynomial
% category is « stricte » when these preliminary conditions allow ⊗ and Γ^*
% to be defined, and an additive P_1 with ⊗ and Γ^* so structured is a
% « ⊗-pd-catégorie »: the two notions are equivalent. Definition of a
% k-linear such category (k -> End(id_{P_1}), with homogeneity of degree
% ν_1...ν_r on the P_{ν_1...ν_r}, p. 84). Example: k-modules; K-Modules on a
% topos, graded K-Modules with Γ^n(M^i) in degree ni. Open questions in the
% « abstract » case: a divided-power structure Γ^p(Γ^q(M)) -> Γ^{pq}(M) on
% Γ^*(M); whether composition P(X,Y) × P(Y,Z) -> P(X,Z) can be seen as a
% linear map P(Y,Z) -> (k-polynomial maps)(P(X,Y), P(X,Z)); reduced to making
% f ↦ Γ^p(f) a polynomial map of degree p between k-modules Hom_1 (p. 85).
%
% Notation fixed once for the batch: Γ^p, Γ^{pq}, α_{p,q} and α^X_{q,p} as on
% p. 80; P_1 (his script P, set \mathcal{P}_1), P_p(X,Y), P_{ν_1...ν_r};
% \mathbf{1} for his double-stroked unit; « ACU » and « pd » kept as written.
% On pp. 81-82 several indices retraced in blue look like « M » where the
% argument requires p (α_{p,q}, P_p); they are read p and noted once.
%
% Readings to check first: « ACU » and « ACU strict » on p. 81; the underlined
% interlinear phrase under Γ^i(0); « un iso » over α_{q,0}; the first words
% of p. 82 and the whole paragraph after the diagram (fast pencil); the phrase
% « stricts que ... » on p. 83; the circled D and the struck lines of p. 84;
% « (M ∈ Ob P_1) » on p. 84; « suppl. » (?) on p. 85.

\documentclass[11pt,a4paper]{article}
% Le préambule commun à toutes les transcriptions du fonds Grothendieck.
%
% Il tient en une idée : **l'incertitude se compose, elle ne se résout pas.**
% Une transcription qui lisse un mot douteux a détruit la seule chose pour
% laquelle on la faisait — un lecteur doit pouvoir savoir, à chaque endroit,
% ce qui était sur le feuillet et ce qui a été ajouté. Les six macros qui
% suivent sont donc l'appareil critique, et non de la décoration.
%
% Les mêmes commandes sont reconnues par `scripts/render.mjs`, qui produit la
% vue de lecture du site. Une macro ajoutée ici doit l'être aussi là-bas,
% faute de quoi le rendu échouera bruyamment — ce qui est voulu : un rendu
% silencieusement faux est pire qu'un rendu absent.

\ProvidesPackage{grothendieck}[2026/08/08 Transcriptions du fonds Grothendieck]

\usepackage{amsmath,amssymb,amsthm}
\usepackage{mathrsfs}
\usepackage{stmaryrd}
% Les diagrammes commutatifs sont partout dans ce fonds : c'est la langue
% même de la géométrie algébrique de Grothendieck, et une transcription qui
% les rendrait en prose ne transcrirait rien.
\usepackage{tikz-cd}
% Français par défaut — c'est la langue des feuillets — mais l'anglais est
% chargé aussi : la traduction et le résumé sont des documents anglais, et la
% typographie française (espaces avant les deux-points) y serait une faute.
% Ces documents basculent par \selectlanguage{english} en tête de corps.
\usepackage[english,french]{babel}
\usepackage{fontspec}
\usepackage{geometry}
\geometry{a4paper,margin=25mm,marginparwidth=28mm}
\usepackage{marginnote}
\usepackage{xcolor}
% ulem, not soul. `soul`'s \st and \ul are built on a character-by-character
% scan that XeLaTeX's font machinery breaks: the document fails with an
% undefined control sequence rather than degrading. ulem does the same two jobs
% and is Unicode-engine safe. `normalem` stops it hijacking \emph, which the
% transcriptions use for Grothendieck's own emphasis.
\usepackage[normalem]{ulem}
\usepackage{hyperref}
\hypersetup{colorlinks=true,linkcolor=black,urlcolor=black,pdfborder={0 0 0}}

% --- Métadonnées du lot -----------------------------------------------------
% Reprises telles quelles par le script de rendu, qui les affiche en tête de
% la vue de lecture. Elles fixent ce que le fichier prétend couvrir : sans
% elles, une transcription flotte sans point d'attache dans un fonds de seize
% mille feuillets.

\newcommand{\folder}[1]{\gdef\grothFolder{#1}}
\newcommand{\batch}[1]{\gdef\grothBatch{#1}}
\newcommand{\leaves}[2]{\gdef\grothFirst{#1}\gdef\grothLast{#2}}
\newcommand{\foldertitle}[1]{\gdef\grothFoldertitle{#1}}
\gdef\grothFolder{?}\gdef\grothBatch{?}\gdef\grothFirst{?}\gdef\grothLast{?}\gdef\grothFoldertitle{}

% --- Le résumé --------------------------------------------------------------
%
% Il ouvre la lecture modernisée, sous le titre « Résumé », et il est composé
% en petit. Ce n'est pas un ornement : c'est le seul passage du document qui
% ne soit pas des mathématiques, et un lecteur doit pouvoir le reconnaître
% avant de l'avoir lu — puis le sauter s'il connaît déjà le sujet, ou s'y
% arrêter s'il ne le connaît pas. Le filet qui le suit marque la reprise du
% corps.

\newenvironment{resume}
  {\small\setlength{\parskip}{0.45em}}
  {\par\medskip\hrule\medskip}

% --- Mots-clés ---------------------------------------------------------------
%
% En anglais, en clôture du résumé de la lecture modernisée : le vocabulaire
% moderne sous lequel chercher ce que les feuillets construisent. Le manifeste
% du site (`npm run manifest`) les extrait pour étiqueter la cote entière —
% cette ligne est la source unique des tags, il n'y a pas de fichier à côté.
\newcommand{\keywords}[1]{\par\smallskip\noindent
  {\footnotesize\textsc{Keywords} --- \textit{#1}}\par}

% --- L'appareil critique ----------------------------------------------------

% Le numéro de feuillet, en marge. C'est l'ancre qui permet de régler un
% désaccord sur une formule en regardant *un* feuillet plutôt qu'une chemise
% de six cents. Le site s'en sert aussi pour tourner les pages du fac-similé
% quand on fait défiler la transcription.
\newcommand{\leaf}[1]{%
  \marginnote{\footnotesize\color{gray!70}\textsf{#1}}%
  \label{leaf:#1}\ignorespaces}

% Illisible : on ne devine pas. Un mot inventé qui se lit comme les autres est
% le pire résultat possible.
\newcommand{\ill}{\textcolor{red!60!black}{[\,\ldots\,]}}

% Lecture douteuse : le mot est donné, et signalé comme incertain.
\newcommand{\uncertain}[1]{\uline{#1}}

% Ajout de l'éditeur — une lettre manquante, un mot sous-entendu.
\newcommand{\add}[1]{\textcolor{blue!50!black}{[#1]}}

% Biffé par Grothendieck lui-même. Ce qu'il a rayé fait partie du document :
% une rature dit souvent où la pensée a changé de direction.
\newcommand{\struck}[1]{\sout{#1}}

% Note du transcripteur, sur l'état matériel du feuillet.
\newcommand{\note}[1]{\par\smallskip{\small\color{gray!60!black}\textit{[#1]}}\par\smallskip}

% Note marginale de Grothendieck, distincte du corps du texte.
\newcommand{\marginal}[1]{\par\smallskip\noindent\hspace{1em}%
  {\small\color{gray!60!black}#1}\par\smallskip}

% --- Titre ------------------------------------------------------------------

\AtBeginDocument{%
  \begin{center}
    {\large\bfseries Fonds Alexandre Grothendieck --- cote n\textsuperscript{o}~\grothFolder}\\[2pt]
    {\itshape\grothFoldertitle}\\[4pt]
    {\small feuillets \grothFirst--\grothLast\ (lot \grothBatch)}\\[2pt]
    {\footnotesize\color{gray!60!black}Transcription établie d'après le fac-similé numérisé
     par l'Université de Montpellier.\\
     Les lectures incertaines sont soulignées, les illisibles notées [\,\ldots\,],
     les ajouts entre crochets.}
  \end{center}
  \vspace{1em}}

\endinput


\folder{137}
\batch{5}
\pages{81}{85}
\dating{[vers 1975-1976]}
\watermark{Édition de démonstration}
\foldertitle{Dévissage des complexes ½ simpliciaux ∂ - parfaits et structures multiplicatives… (1975 ou 1976) : notes manuscrites (s.d.)}

\begin{document}

\page{81}
Inversement, soit $\mathcal{P}_1$ \struck{\ill{}} catégorie $\otimes$-additive
\uncertain{ACU}, munie d'un $\otimes$-foncteur \struck{\ill{}}
\uncertain{ACU strict} de $(\mathcal{P}_1, +)$ dans
$(\mathcal{P}_{1+}^{\mathrm{grad.}}, \otimes)$, \struck{\ill{}}
\note{« ACU » est écrit au-dessus de la ligne, après « additive », avec un
signe d'insertion ; « ACU strict » est porté au-dessus de « foncteur », à la
place d'un mot biffé, le second mot étant lui-même retouché. L'indice $+$ sous
$\mathcal{P}_1^{\mathrm{grad.}}$ est d'une lecture douteuse.}
\[
\Gamma^{*} : (\mathcal{P}_1, \oplus) \longrightarrow
(\mathcal{P}_1^{\mathrm{grad.}}, \otimes)
\]
donc
\[
\Gamma^{*}(0) = \mathbf{1} \quad \text{i.e.} \quad
\Gamma^{i}(0) = \begin{cases} \mathbf{1} & \text{si } i = 0 \\
0 & \text{si } i > 0 \end{cases}
\]
\ill{}
\note{sous la formule, une ligne soulignée, interlinéaire, de trois ou quatre
mots, que l'on n'a pas su lire}
\[
\Gamma^{*}(X \oplus Y) \simeq \Gamma^{*}(X) \otimes \Gamma^{*}(Y)
\]
(isom.\ \uncertain{gr.}\ \ill{}, fonctoriel en $X, Y$, compatible avec
contraintes d'unité, d'ass., de comm.) avec
\[
\Gamma^{i}(M) = \begin{cases} 0 & \text{si } i < 0 \\
\mathbf{1} & \text{si } i = 0 \\
M & \text{si } i = 1 \end{cases}
\]
(isom.\ de foncteurs en $M$). Supposons de plus donnés des homs de foncteurs
\[
\Gamma^{pq} \xrightarrow{\ \alpha_{pq}\ } \Gamma^{p} \circ \Gamma^{q}
\]
(i.e.
\[
\Gamma^{pq}(X) \xrightarrow{\ \alpha^{X}_{pq}\ } \Gamma^{p}(\Gamma^{q}(X))\ )
\]
\note{l'indice de $\alpha$, repassé au stylo bleu, a la forme d'un « M » ; la
page 80 écrit $\alpha^{X}_{p,q}$ et c'est ce que demande le contexte. Le but
$\Gamma^{p} \circ \Gamma^{q}$ est corrigé, sur la page, depuis
$\Gamma^{p}(\Gamma^{q}$.}
satisfaisant à la condition d'associativité ci-dessus, plus ces conditions
d'unité pour \struck{$\Gamma^{pq}$} l'un des indices $p, q$ égaux à $0$ ou
$1$ :
\note{« la condition d'associativité ci-dessus » est la condition (*) de la
page 80}
\begin{align*}
(**) \qquad
& \Gamma^{q} \xrightarrow{\ \alpha_{1,q} = \mathrm{id}\ }
\Gamma^{1} \circ \Gamma^{q} = \mathrm{id} \circ \Gamma^{q} = \Gamma^{q} \\
& \Gamma^{q} \xrightarrow{\ \alpha_{q,1} = \mathrm{id}\ }
\Gamma^{q} \circ \Gamma^{1} = \Gamma^{q} \circ \mathrm{id} = \Gamma^{q} \\[1ex]
(***) \qquad
& \Gamma^{0} \xrightarrow{\ \alpha_{0,q} = \mathrm{id}\ }
\Gamma^{0} \circ \Gamma^{q} = \Gamma^{0} \\
& \Gamma^{0} \xrightarrow[\ \simeq\ ]{\ \alpha_{q,0}\ }
\Gamma^{q} \circ \Gamma^{0} .
\end{align*}
\note{dans chacune des deux conditions, la seconde flèche part de la même
source que la première et s'en détache par une courbe. Sur la dernière,
après $\alpha_{q,0}$, deux petits mots au-dessus de la flèche, lus
\uncertain{un iso}, et un signe $\simeq$ souligné deux fois.}

\page{82}
Les $\alpha_{p,q}$ permettent de définir des \ill{}
\note{l'indice $p$ de $\alpha_{p,q}$ est retouché au stylo bleu ; la ligne
suivante, courte et soulignée, commence par un mot surchargé, peut-être
« Définissons », suivi d'un « p » (ou $\mathcal{P}$) de lecture incertaine}
\[
\mathcal{P}_p(X, Y) = \mathrm{Hom}(\Gamma^{p} X, Y)
\]
\note{l'indice $p$ de $\mathcal{P}_p$ est repassé en bleu et a la forme d'un
« M », comme à la page 81}
donc
\[
\begin{cases}
\mathcal{P}_1(X, Y) = \mathrm{Hom}_{\mathcal{P}_1}(X, Y) \\
\mathcal{P}_0(X, Y) = \mathrm{Hom}_{\mathcal{P}_1}(\mathbf{1}, Y)
\end{cases}
\]
\note{après chacune des deux égalités, un petit groupe de signes illisible,
à droite ; le second est souligné d'une accolade}
et des applications \uncertain{bil.}
\[
\begin{array}{ccc}
\mathcal{P}_p(X, Y) \times \mathcal{P}_q(Y, Z) & \longrightarrow &
\mathcal{P}_{pq}(X, Z) \\
(f, g) & \longmapsto & g \circ f
\end{array}
\]
par
\[
\Gamma^{p} X \xrightarrow{\ f\ } Y \qquad \Gamma^{q} Y \xrightarrow{\ g\ } Z
\]

\begin{tikzcd}
 & \Gamma^{q} Y \arrow[r, "g"] & Z \\
\Gamma^{q}(\Gamma^{p}(X)) \arrow[ur, "\Gamma^{q}(f)"] & & \\
\Gamma^{pq}(X) \arrow[u, "\alpha^{X}_{q,p}"'] & &
\end{tikzcd}

\[
g \circ f = g\, \Gamma^{q}(f)\, \alpha^{X}_{q,p}
\]
\note{la page trace d'abord, à droite de $\Gamma^{q}(\Gamma^{p}(X))$, une flèche
horizontale vers $\Gamma^{q}(Y)$ suivie d'une flèche vers $Z$, puis la biffe
au stylo bleu et la remplace par la flèche oblique $\Gamma^{q}(f)$ ; on a
dessiné le diagramme corrigé. L'exposant de $\Gamma^{pq}(X)$ est retouché.}

Le diagramme commute (*) assure l'associativité de la composition, (**)
\uncertain{signifie} \ill{} que si \uncertain{$p = 1$ ou $q = 1$}
\struck{\ill{}} \struck{\ill{}}, dans le composé \uncertain{est celui}
provenant de la bifonctorialité de $\mathcal{P}_p(X, Y)$ sur
$\mathcal{P}_1^{\circ} \times \mathcal{P}_1$, enfin la compatibilité (***)
signifie, \ill{} ce \uncertain{dernier}, que
\[
g \circ f = g \quad \text{si } g \text{ de degré } 0
\]
\note{le paragraphe est au crayon, d'une écriture rapide ; « $p = 1$ ou
$q = 1$ » est interlinéaire, au stylo bleu, au-dessus d'un passage biffé}
la deuxième signifie
\[
\mathbf{1} \xrightarrow[\ \simeq\ ]{\ \alpha_{q,0}\ } \Gamma^{q}(\mathbf{1})
\]
\note{la formule est entourée d'un grand ovale au crayon ; au-dessus, un
indice biffé}
\marginal{pour \ill{}}
je l'avais oublié dans les axiomes préliminaires, il me semble, pour \ill{}
\note{le renvoi en marge gauche, écrit de biais, est relié au texte par deux
traits verticaux}
\marginal{\uncertain{pour la} \ill{} des \ill{}}
\note{cette dernière ligne est d'une autre encre, noire, et d'un autre trait
que le reste de la page}

\page{83}
Enfin, si $X = \bigoplus_{1 \leq i \leq r} X_i$, la décomposition de
$\Gamma^{n}(X)$ en \struck{les} $\bigotimes \Gamma^{\nu_i}(X_i)$ donne
\struck{une} des décompositions de $\mathcal{P}_n(\prod X_i, Y)$ en des
$\mathcal{P}_{\nu_1 \ldots \nu_r}(X_1, \ldots, X_r, Y)$.

On vérifie alors les axiomes a) à f), et le fait qu'\struck{il y a}
les produits tensoriels de $\mathcal{P}_1$ \uncertain{sont} \emph{stricts}
\ill{} (\ill{} \uncertain{foncteurs}), que les $\Gamma^{\nu_i}$ existent et
satisfont la condition énoncée précédemment \ldots
\note{les axiomes a) à f) ne sont pas sur ces feuillets}

Une catégorie polynomiale $\mathcal{P}$ est dite \emph{stricte} si elle
satisfait aux conditions préliminaires permettant de définir $\otimes$ et
$\Gamma^{*}$ avec toutes les propriétés envisagées. D'autre part une
catégorie additive $\mathcal{P}_1$ munie de $\otimes$ et $\Gamma^{*}$ avec les
structures décrites, est dite une $\otimes$-pd-catégorie. Les deux notions
sont donc équivalentes.

On dit que $\mathcal{P}$ est \struck{\ill{}} une catégorie
\uncertain{polynomiale} (resp.\ $\otimes$-pd-cat.)
\uncertain{$k$-linéaire} ($k$ anneau commutatif) si on s'est donné un hom
d'anneaux
\[
k \to \mathrm{End}(\mathrm{id}_{\mathcal{P}_1}) \quad
(\text{ou } \mathrm{id}_{\mathcal{P}})
\]

\page{84}
tel que la décomposition de
$\mathcal{P}_n(X_1 \times \cdots \times X_r; Y)$ en les
$\mathcal{P}_{\nu_1 \ldots \nu_r}(X_1, \ldots, X_r; Y)$ respecte les op.\ de
$k$, i.e.\ $f \in \mathcal{P}_{\nu_1 \ldots \nu_r}$ implique
\[
f \circ (\mathrm{id}_{\lambda_1} \times \cdots \times \mathrm{id}_{\lambda_r})
= \lambda_1^{\nu_1} \cdots \lambda_r^{\nu_r} f
\]
pour tout $(\lambda_1, \ldots, \lambda_r) \in k^{r}$.
\note{$\mathrm{id}_{\lambda_i}$ comme sur la page : l'homothétie de rapport
$\lambda_i$}

\emph{Ex.} La catégorie des $k$-modules est une $\otimes$-pd-catégorie.

(D) $\Rightarrow$ si $X$ est un topos avec une \uncertain{$k$-Algèbre} $K$, la
cat.\ des $K$-Modules
\note{« (D) » : un D majuscule cerclé, suivi d'une double flèche ; le mot
« $k$-Algèbre » est surchargé}
\struck{\ill{}} [(D) $\Rightarrow$ \uncertain{pour} prendre des $K$-Modules
gradués, \struck{par \ill{}}
\marginal{?}
\struck{quelconques \ill{}}, on considère des \ill{} dans
$\Gamma^{n}(M^{i})$ pour le degré $ni$.]

Il y a diverses questions dans le cas « abstrait ». Par exemple : quand
\uncertain{$(M \in \mathrm{Ob}\,\mathcal{P}_1)$} \uncertain{peut-on}
\ill{} sur $\Gamma^{*}(M)$ une structure « à p.d. », i.e.\ un hom
\uncertain{de degré} \ill{}
\[
\Gamma^{p}(\Gamma^{q}(M)) \longrightarrow \Gamma^{pq}(M)
\]
satisfaisant aux axiomes des p.d.\,?

\page{85}
Y a-t-il une telle structure canonique dans le cas de l'exemple ou
\uncertain{abstrait} : quand peut-on \uncertain{considérer} la composition
des morphismes de $\mathcal{P}$
\[
\mathcal{P}(X, Y) \times \mathcal{P}(Y, Z) \longrightarrow \mathcal{P}(X, Z)
\]
\note{la ligne commence par un « $\mathrm{Hom}_{\mathcal{P}}$ » biffé}
($k$-linéaire en la deuxième variable) en une application linéaire
\[
\mathcal{P}(Y, Z) \longrightarrow
(\text{Appl.\ } k\text{-polynomiales})(\mathcal{P}(X, Y), \mathcal{P}(X, Z))
\]
(Ceci \ill{} acquis dans l'exemple principal envisagé). La construction
devrait se ramener pour l'essentiel à définir de telles structures
\uncertain{suppl.}\ sur les foncteurs $\Gamma^{n}$, i.e.\ pour tout $X, Y$
\[
f \longmapsto \Gamma^{p}(f) : \mathrm{Hom}_1(X, Y) \longrightarrow
\mathrm{Hom}_1(\Gamma^{p} X, \Gamma^{p} Y)
\]
doit être défini comme « application polynomiale de degré $p$ » entre les
$k$-modules $\mathrm{Hom}_1(X, Y)$ et
$\mathrm{Hom}_1(\Gamma^{p} X, \Gamma^{p} Y)$, et non seulement comme
application ensembliste.
\note{la page s'arrête ici, au premier tiers ; c'est la dernière du dossier}

\end{document}
